% Options for packages loaded elsewhere
\PassOptionsToPackage{unicode}{hyperref}
\PassOptionsToPackage{hyphens}{url}
\PassOptionsToPackage{dvipsnames,svgnames*,x11names*}{xcolor}
%
\documentclass[
  a4paper,
  UTF8]{ctexart}
\usepackage{amsmath,amssymb}
\usepackage{lmodern}
\usepackage{iftex}
\ifPDFTeX
  \usepackage[T1]{fontenc}
  \usepackage[utf8]{inputenc}
  \usepackage{textcomp} % provide euro and other symbols
\else % if luatex or xetex
  \usepackage{unicode-math}
  \defaultfontfeatures{Scale=MatchLowercase}
  \defaultfontfeatures[\rmfamily]{Ligatures=TeX,Scale=1}
\fi
% Use upquote if available, for straight quotes in verbatim environments
\IfFileExists{upquote.sty}{\usepackage{upquote}}{}
\IfFileExists{microtype.sty}{% use microtype if available
  \usepackage[]{microtype}
  \UseMicrotypeSet[protrusion]{basicmath} % disable protrusion for tt fonts
}{}
\makeatletter
\@ifundefined{KOMAClassName}{% if non-KOMA class
  \IfFileExists{parskip.sty}{%
    \usepackage{parskip}
  }{% else
    \setlength{\parindent}{0pt}
    \setlength{\parskip}{6pt plus 2pt minus 1pt}}
}{% if KOMA class
  \KOMAoptions{parskip=half}}
\makeatother
\usepackage{xcolor}
\IfFileExists{xurl.sty}{\usepackage{xurl}}{} % add URL line breaks if available
\IfFileExists{bookmark.sty}{\usepackage{bookmark}}{\usepackage{hyperref}}
\hypersetup{
  pdftitle={接口文档},
  pdfauthor={DBMS 控制台项目},
  colorlinks=true,
  linkcolor={blue},
  filecolor={Maroon},
  citecolor={Blue},
  urlcolor={blue},
  pdfcreator={LaTeX via pandoc}}
\urlstyle{same} % disable monospaced font for URLs
\usepackage[margin=2.5cm]{geometry}
\usepackage{color}
\usepackage{fancyvrb}
\newcommand{\VerbBar}{|}
\newcommand{\VERB}{\Verb[commandchars=\\\{\}]}
\DefineVerbatimEnvironment{Highlighting}{Verbatim}{commandchars=\\\{\}}
% Add ',fontsize=\small' for more characters per line
\newenvironment{Shaded}{}{}
\newcommand{\AlertTok}[1]{\textcolor[rgb]{1.00,0.00,0.00}{\textbf{#1}}}
\newcommand{\AnnotationTok}[1]{\textcolor[rgb]{0.38,0.63,0.69}{\textbf{\textit{#1}}}}
\newcommand{\AttributeTok}[1]{\textcolor[rgb]{0.49,0.56,0.16}{#1}}
\newcommand{\BaseNTok}[1]{\textcolor[rgb]{0.25,0.63,0.44}{#1}}
\newcommand{\BuiltInTok}[1]{#1}
\newcommand{\CharTok}[1]{\textcolor[rgb]{0.25,0.44,0.63}{#1}}
\newcommand{\CommentTok}[1]{\textcolor[rgb]{0.38,0.63,0.69}{\textit{#1}}}
\newcommand{\CommentVarTok}[1]{\textcolor[rgb]{0.38,0.63,0.69}{\textbf{\textit{#1}}}}
\newcommand{\ConstantTok}[1]{\textcolor[rgb]{0.53,0.00,0.00}{#1}}
\newcommand{\ControlFlowTok}[1]{\textcolor[rgb]{0.00,0.44,0.13}{\textbf{#1}}}
\newcommand{\DataTypeTok}[1]{\textcolor[rgb]{0.56,0.13,0.00}{#1}}
\newcommand{\DecValTok}[1]{\textcolor[rgb]{0.25,0.63,0.44}{#1}}
\newcommand{\DocumentationTok}[1]{\textcolor[rgb]{0.73,0.13,0.13}{\textit{#1}}}
\newcommand{\ErrorTok}[1]{\textcolor[rgb]{1.00,0.00,0.00}{\textbf{#1}}}
\newcommand{\ExtensionTok}[1]{#1}
\newcommand{\FloatTok}[1]{\textcolor[rgb]{0.25,0.63,0.44}{#1}}
\newcommand{\FunctionTok}[1]{\textcolor[rgb]{0.02,0.16,0.49}{#1}}
\newcommand{\ImportTok}[1]{#1}
\newcommand{\InformationTok}[1]{\textcolor[rgb]{0.38,0.63,0.69}{\textbf{\textit{#1}}}}
\newcommand{\KeywordTok}[1]{\textcolor[rgb]{0.00,0.44,0.13}{\textbf{#1}}}
\newcommand{\NormalTok}[1]{#1}
\newcommand{\OperatorTok}[1]{\textcolor[rgb]{0.40,0.40,0.40}{#1}}
\newcommand{\OtherTok}[1]{\textcolor[rgb]{0.00,0.44,0.13}{#1}}
\newcommand{\PreprocessorTok}[1]{\textcolor[rgb]{0.74,0.48,0.00}{#1}}
\newcommand{\RegionMarkerTok}[1]{#1}
\newcommand{\SpecialCharTok}[1]{\textcolor[rgb]{0.25,0.44,0.63}{#1}}
\newcommand{\SpecialStringTok}[1]{\textcolor[rgb]{0.73,0.40,0.53}{#1}}
\newcommand{\StringTok}[1]{\textcolor[rgb]{0.25,0.44,0.63}{#1}}
\newcommand{\VariableTok}[1]{\textcolor[rgb]{0.10,0.09,0.49}{#1}}
\newcommand{\VerbatimStringTok}[1]{\textcolor[rgb]{0.25,0.44,0.63}{#1}}
\newcommand{\WarningTok}[1]{\textcolor[rgb]{0.38,0.63,0.69}{\textbf{\textit{#1}}}}
\usepackage{longtable,booktabs,array}
\usepackage{calc} % for calculating minipage widths
% Correct order of tables after \paragraph or \subparagraph
\usepackage{etoolbox}
\makeatletter
\patchcmd\longtable{\par}{\if@noskipsec\mbox{}\fi\par}{}{}
\makeatother
% Allow footnotes in longtable head/foot
\IfFileExists{footnotehyper.sty}{\usepackage{footnotehyper}}{\usepackage{footnote}}
\makesavenoteenv{longtable}
\setlength{\emergencystretch}{3em} % prevent overfull lines
\providecommand{\tightlist}{%
  \setlength{\itemsep}{0pt}\setlength{\parskip}{0pt}}
\setcounter{secnumdepth}{5}
\ifLuaTeX
  \usepackage{selnolig}  % disable illegal ligatures
\fi

\title{接口文档}
\author{DBMS 控制台项目}
\date{}

\begin{document}
\maketitle

\renewcommand*\contentsname{目录}
{
\hypersetup{linkcolor=}
\setcounter{tocdepth}{3}
\tableofcontents
}
\hypertarget{ux63a5ux53e3ux89c4ux8303}{%
\section{接口规范}\label{ux63a5ux53e3ux89c4ux8303}}

本文档是项目唯一权威接口文档。早期登录注册草稿《接口文档.md》已合并到本文档，不再单独维护，避免接口设计和接口文档重复导致实现偏差。

\hypertarget{ux63a5ux53e3ux603bux5219}{%
\subsection{1. 接口总则}\label{ux63a5ux53e3ux603bux5219}}

\hypertarget{base-url}{%
\subsubsection{1.1 Base URL}\label{base-url}}

\begin{Shaded}
\begin{Highlighting}[]
\NormalTok{/api/v1}
\end{Highlighting}
\end{Shaded}

\hypertarget{ux6210ux529fux54cdux5e94}{%
\subsubsection{1.2 成功响应}\label{ux6210ux529fux54cdux5e94}}

\begin{Shaded}
\begin{Highlighting}[]
\FunctionTok{\{}
  \DataTypeTok{"code"}\FunctionTok{:} \DecValTok{200}\FunctionTok{,}
  \DataTypeTok{"message"}\FunctionTok{:} \StringTok{"success"}\FunctionTok{,}
  \DataTypeTok{"data"}\FunctionTok{:} \FunctionTok{\{\},}
  \DataTypeTok{"traceId"}\FunctionTok{:} \StringTok{"01HZ0000000000000000000000"}
\FunctionTok{\}}
\end{Highlighting}
\end{Shaded}

\begin{longtable}[]{@{}llll@{}}
\toprule
字段 & 类型 & 必填 & 约束 \\
\midrule
\endhead
code & number & 是 & 成功固定为 200 \\
message & string & 是 & 成功提示 \\
data & any & 是 & 无数据时为 null \\
traceId & string & 是 & 服务端生成，请求级唯一 \\
\bottomrule
\end{longtable}

\hypertarget{ux9519ux8befux54cdux5e94}{%
\subsubsection{1.3 错误响应}\label{ux9519ux8befux54cdux5e94}}

\begin{Shaded}
\begin{Highlighting}[]
\FunctionTok{\{}
  \DataTypeTok{"code"}\FunctionTok{:} \DecValTok{400}\FunctionTok{,}
  \DataTypeTok{"message"}\FunctionTok{:} \StringTok{"字段 age 不允许为空"}\FunctionTok{,}
  \DataTypeTok{"data"}\FunctionTok{:} \KeywordTok{null}\FunctionTok{,}
  \DataTypeTok{"traceId"}\FunctionTok{:} \StringTok{"01HZ0000000000000000000000"}\FunctionTok{,}
  \DataTypeTok{"error"}\FunctionTok{:} \FunctionTok{\{}
    \DataTypeTok{"type"}\FunctionTok{:} \StringTok{"MYSQL\_CONSTRAINT\_ERROR"}\FunctionTok{,}
    \DataTypeTok{"code"}\FunctionTok{:} \StringTok{"DB\_NOT\_NULL\_VIOLATION"}\FunctionTok{,}
    \DataTypeTok{"source"}\FunctionTok{:} \StringTok{"mysql"}\FunctionTok{,}
    \DataTypeTok{"details"}\FunctionTok{:} \StringTok{"Column \textquotesingle{}age\textquotesingle{} cannot be null"}\FunctionTok{,}
    \DataTypeTok{"fieldErrors"}\FunctionTok{:} \OtherTok{[}
      \FunctionTok{\{}
        \DataTypeTok{"field"}\FunctionTok{:} \StringTok{"age"}\FunctionTok{,}
        \DataTypeTok{"message"}\FunctionTok{:} \StringTok{"字段 age 不允许为空"}
      \FunctionTok{\}}
    \OtherTok{]}\FunctionTok{,}
    \DataTypeTok{"mysql"}\FunctionTok{:} \FunctionTok{\{}
      \DataTypeTok{"errno"}\FunctionTok{:} \DecValTok{1048}\FunctionTok{,}
      \DataTypeTok{"sqlState"}\FunctionTok{:} \StringTok{"23000"}\FunctionTok{,}
      \DataTypeTok{"sqlMessage"}\FunctionTok{:} \StringTok{"Column \textquotesingle{}age\textquotesingle{} cannot be null"}\FunctionTok{,}
      \DataTypeTok{"databaseName"}\FunctionTok{:} \StringTok{"student"}\FunctionTok{,}
      \DataTypeTok{"tableName"}\FunctionTok{:} \StringTok{"students"}\FunctionTok{,}
      \DataTypeTok{"columnName"}\FunctionTok{:} \StringTok{"age"}\FunctionTok{,}
      \DataTypeTok{"constraintName"}\FunctionTok{:} \KeywordTok{null}
    \FunctionTok{\}}
  \FunctionTok{\}}
\FunctionTok{\}}
\end{Highlighting}
\end{Shaded}

\begin{longtable}[]{@{}llll@{}}
\toprule
字段 & 类型 & 必填 & 约束 \\
\midrule
\endhead
code & number & 是 & 与 HTTP 状态语义一致 \\
message & string & 是 & 前端直接展示的精准提示 \\
data & null & 是 & 错误时固定为 null \\
traceId & string & 是 & 请求级唯一 \\
error.type & string & 是 & 错误类型枚举 \\
error.code & string & 是 & 错误码枚举 \\
error.source & string & 是 & client、server、mysql \\
error.details & string & 是 & 错误详情，不包含密码和 Token \\
error.fieldErrors & array & 是 & 字段级错误列表 \\
error.mysql & object & 是 & 非 MySQL 错误时为 null \\
\bottomrule
\end{longtable}

\hypertarget{ux524dux7aefux672cux5730ux9519ux8befux6a21ux578b}{%
\subsubsection{1.4
前端本地错误模型}\label{ux524dux7aefux672cux5730ux9519ux8befux6a21ux578b}}

以下错误没有后端响应，由前端统一构造展示：

\begin{itemize}
\tightlist
\item
  无法连接后端：\texttt{type=NETWORK\_ERROR}；\texttt{code=CLIENT\_BACKEND\_UNREACHABLE}；\texttt{message=无法连接后端服务，请检查服务是否启动}。
\item
  请求超时：\texttt{type=HTTP\_TIMEOUT}；\texttt{code=CLIENT\_REQUEST\_TIMEOUT}；\texttt{message=后端响应超时，请稍后重试}。
\item
  请求被浏览器取消：\texttt{type=CLIENT\_ABORTED}；\texttt{code=CLIENT\_REQUEST\_ABORTED}；\texttt{message=请求已取消}。
\end{itemize}

\hypertarget{http-ux72b6ux6001ux7801}{%
\subsubsection{1.5 HTTP 状态码}\label{http-ux72b6ux6001ux7801}}

\begin{longtable}[]{@{}lll@{}}
\toprule
HTTP & error.type & 使用场景 \\
\midrule
\endhead
400 & VALIDATION\_ERROR & 参数格式错误、字段校验失败 \\
401 & AUTH\_ERROR & 未登录、Access Token 失效、登录凭证错误 \\
403 & AUTHZ\_ERROR & 已登录但无权访问资源 \\
404 & RESOURCE\_NOT\_FOUND & 数据库、表、字段、索引、视图不存在 \\
409 & RESOURCE\_CONFLICT & 名称重复、状态冲突 \\
412 & CONFIRMATION\_REQUIRED & 缺少确认信息 \\
413 & LIMIT\_EXCEEDED & 容量超限 \\
429 & RATE\_LIMITED & 请求过于频繁 \\
500 & INTERNAL\_ERROR & 服务端未预期异常 \\
502 & MYSQL\_CONNECTION\_ERROR & MySQL 连接失败 \\
503 & MYSQL\_UNAVAILABLE & MySQL 不可用 \\
\bottomrule
\end{longtable}

\hypertarget{ux901aux7528ux53c2ux6570ux7ea6ux675f}{%
\subsection{2.
通用参数约束}\label{ux901aux7528ux53c2ux6570ux7ea6ux675f}}

\hypertarget{ux6807ux8bc6ux7b26}{%
\subsubsection{2.1 标识符}\label{ux6807ux8bc6ux7b26}}

数据库显示名、表名、字段名、索引名、约束名、视图名：

\begin{Shaded}
\begin{Highlighting}[]
\NormalTok{\^{}[A{-}Za{-}z][A{-}Za{-}z0{-}9\_]\{0,63\}$}
\end{Highlighting}
\end{Shaded}

\hypertarget{ux5206ux9875}{%
\subsubsection{2.2 分页}\label{ux5206ux9875}}

\begin{longtable}[]{@{}lllll@{}}
\toprule
参数 & 类型 & 必填 & 默认值 & 约束 \\
\midrule
\endhead
page & number & 否 & 1 & 整数，最小 1 \\
pageSize & number & 否 & 20 & 整数，最小 1，最大 100 \\
\bottomrule
\end{longtable}

\hypertarget{ux786eux8ba4ux5bf9ux8c61}{%
\subsubsection{2.3 确认对象}\label{ux786eux8ba4ux5bf9ux8c61}}

高风险请求统一使用：

\begin{Shaded}
\begin{Highlighting}[]
\FunctionTok{\{}
  \DataTypeTok{"confirmation"}\FunctionTok{:} \FunctionTok{\{}
    \DataTypeTok{"confirmed"}\FunctionTok{:} \KeywordTok{true}\FunctionTok{,}
    \DataTypeTok{"confirmText"}\FunctionTok{:} \StringTok{"students"}\FunctionTok{,}
    \DataTypeTok{"disableSameTypePrompt"}\FunctionTok{:} \KeywordTok{false}
  \FunctionTok{\}}
\FunctionTok{\}}
\end{Highlighting}
\end{Shaded}

\begin{longtable}[]{@{}llll@{}}
\toprule
字段 & 类型 & 必填 & 约束 \\
\midrule
\endhead
confirmed & boolean & 是 & 必须为 true \\
confirmText & string & 是 & 必须与后端要求文本完全一致 \\
disableSameTypePrompt & boolean & 否 & 只对中风险操作生效 \\
\bottomrule
\end{longtable}

强制确认操作忽略 \texttt{disableSameTypePrompt}。

\hypertarget{ux9ed8ux8ba4ux5bb9ux91cfux9650ux5236}{%
\subsubsection{2.4
默认容量限制}\label{ux9ed8ux8ba4ux5bb9ux91cfux9650ux5236}}

\begin{longtable}[]{@{}ll@{}}
\toprule
限制项 & 默认值 \\
\midrule
\endhead
单用户数据库数量 & 10 \\
单数据库表数量 & 100 \\
单数据库视图数量 & 50 \\
单表字段数量 & 80 \\
单表索引数量 & 16 \\
单用户物理数据总量 & 1GB \\
单数据库物理数据量 & 512MB \\
单表物理数据量 & 256MB \\
单次分页返回行数 & 100 \\
单次查询最大返回行数 & 1000 \\
单次批量插入行数 & 500 \\
单次批量修改影响行数 & 200 \\
单次批量删除影响行数 & 200 \\
查询最大 JOIN 数 & 8 \\
查询最大子查询深度 & 3 \\
查询最大执行时间 & 10 秒 \\
查询 AST 最大体积 & 64KB \\
头像文件大小 & 2MB \\
背景图文件大小 & 8MB \\
单用户资源总量 & 100MB \\
\bottomrule
\end{longtable}

服务端可以通过配置调整默认值。接口文档中的默认值用于本地课程项目落地。

\hypertarget{ux9519ux8befux7801ux679aux4e3e}{%
\subsection{3. 错误码枚举}\label{ux9519ux8befux7801ux679aux4e3e}}

\hypertarget{ux901aux7528ux9519ux8befux7801}{%
\subsubsection{3.1 通用错误码}\label{ux901aux7528ux9519ux8befux7801}}

\begin{longtable}[]{@{}ll@{}}
\toprule
code & message 模板 \\
\midrule
\endhead
VALIDATION\_FIELD\_REQUIRED & \texttt{\{field\}\ 为必填项} \\
VALIDATION\_IDENTIFIER\_INVALID & \texttt{\{field\}\ 名称格式不正确} \\
VALIDATION\_ENUM\_INVALID & \texttt{\{field\}\ 不在允许范围内} \\
RESOURCE\_NOT\_FOUND & \texttt{\{resource\}\ 不存在} \\
RESOURCE\_CONFLICT & \texttt{\{resource\}\ 已存在} \\
AUTH\_TOKEN\_MISSING & \texttt{登录状态缺失，请重新登录} \\
AUTH\_TOKEN\_EXPIRED & \texttt{登录已过期，请重新登录} \\
AUTH\_CREDENTIAL\_INVALID & \texttt{账号与密码不匹配} \\
AUTHZ\_DATABASE\_FORBIDDEN & \texttt{无权访问该数据库} \\
CONFIRMATION\_TEXT\_MISMATCH & \texttt{确认文本不匹配} \\
LIMIT\_EXCEEDED & \texttt{操作超过系统限制} \\
INTERNAL\_ERROR & \texttt{服务端异常，请联系维护者并提供追踪编号} \\
\bottomrule
\end{longtable}

\hypertarget{mysql-ux9519ux8befux7801ux6620ux5c04}{%
\subsubsection{3.2 MySQL
错误码映射}\label{mysql-ux9519ux8befux7801ux6620ux5c04}}

\begin{itemize}
\tightlist
\item
  \texttt{1044} /
  \texttt{42000}：\texttt{DB\_ACCESS\_DENIED}；\texttt{数据库访问被拒绝}。
\item
  \texttt{1045} /
  \texttt{28000}：\texttt{DB\_AUTH\_FAILED}；\texttt{数据库账号认证失败}。
\item
  \texttt{1046} /
  \texttt{3D000}：\texttt{DB\_NOT\_SELECTED}；\texttt{未选择数据库}。
\item
  \texttt{1048} /
  \texttt{23000}：\texttt{DB\_NOT\_NULL\_VIOLATION}；\texttt{字段\ \{columnName\}\ 不允许为空}。
\item
  \texttt{1050} /
  \texttt{42S01}：\texttt{DB\_TABLE\_EXISTS}；\texttt{表\ \{tableName\}\ 已存在}。
\item
  \texttt{1051} /
  \texttt{42S02}：\texttt{DB\_TABLE\_NOT\_FOUND}；\texttt{表\ \{tableName\}\ 不存在}。
\item
  \texttt{1054} /
  \texttt{42S22}：\texttt{DB\_COLUMN\_NOT\_FOUND}；\texttt{字段\ \{columnName\}\ 不存在}。
\item
  \texttt{1062} /
  \texttt{23000}：\texttt{DB\_DUPLICATE\_ENTRY}；\texttt{唯一约束冲突：\{constraintName\}}。
\item
  \texttt{1064} /
  \texttt{42000}：\texttt{DB\_SQL\_SYNTAX\_ERROR}；\texttt{系统生成的\ SQL\ 语法错误}。
\item
  \texttt{1072} /
  \texttt{42000}：\texttt{DB\_KEY\_COLUMN\_NOT\_FOUND}；\texttt{索引字段\ \{columnName\}\ 不存在}。
\item
  \texttt{1215} /
  \texttt{HY000}：\texttt{DB\_FOREIGN\_KEY\_CREATE\_FAILED}；\texttt{外键创建失败}。
\item
  \texttt{1216} /
  \texttt{23000}：\texttt{DB\_FOREIGN\_KEY\_CHILD\_FAILED}；\texttt{外键约束失败，引用数据不存在}。
\item
  \texttt{1217} /
  \texttt{23000}：\texttt{DB\_FOREIGN\_KEY\_PARENT\_FAILED}；\texttt{外键约束失败，存在关联数据}。
\item
  \texttt{1451} /
  \texttt{23000}：\texttt{DB\_FOREIGN\_KEY\_RESTRICT}；\texttt{存在关联数据，删除被外键限制}。
\item
  \texttt{1452} /
  \texttt{23000}：\texttt{DB\_FOREIGN\_KEY\_REFERENCE\_MISSING}；\texttt{外键引用的目标数据不存在}。
\item
  \texttt{3819} /
  \texttt{HY000}：\texttt{DB\_CHECK\_CONSTRAINT\_FAILED}；\texttt{检查约束\ \{constraintName\}\ 不满足}。
\end{itemize}

\hypertarget{ux8ba4ux8bc1ux63a5ux53e3}{%
\subsection{4. 认证接口}\label{ux8ba4ux8bc1ux63a5ux53e3}}

\hypertarget{api-auth-001-ux6ce8ux518c}{%
\subsubsection{API-AUTH-001 注册}\label{api-auth-001-ux6ce8ux518c}}

\begin{Shaded}
\begin{Highlighting}[]
\NormalTok{POST /api/v1/auth/register}
\end{Highlighting}
\end{Shaded}

请求体：

\begin{longtable}[]{@{}llll@{}}
\toprule
参数 & 类型 & 必填 & 约束 \\
\midrule
\endhead
userName & string & 是 & trim 后长度 2-20 \\
password & string & 是 & 8-16 位，包含英文字母和数字 \\
\bottomrule
\end{longtable}

成功响应 \texttt{data}：

\begin{longtable}[]{@{}lll@{}}
\toprule
字段 & 类型 & 约束 \\
\midrule
\endhead
accountNo & string & 10 位数字 \\
userName & string & 展示昵称 \\
\bottomrule
\end{longtable}

错误：

\begin{longtable}[]{@{}ll@{}}
\toprule
code & message \\
\midrule
\endhead
VALIDATION\_FIELD\_REQUIRED & 昵称和密码均为必填项 \\
VALIDATION\_PASSWORD\_INVALID & 密码必须为 8-16
位且包含英文字母和数字 \\
INTERNAL\_ERROR & 注册失败，请提供追踪编号联系维护者 \\
\bottomrule
\end{longtable}

\hypertarget{api-auth-002-ux767bux5f55}{%
\subsubsection{API-AUTH-002 登录}\label{api-auth-002-ux767bux5f55}}

\begin{Shaded}
\begin{Highlighting}[]
\NormalTok{POST /api/v1/auth/login}
\end{Highlighting}
\end{Shaded}

请求体：

\begin{longtable}[]{@{}llll@{}}
\toprule
参数 & 类型 & 必填 & 约束 \\
\midrule
\endhead
accountNo & string & 是 & 10 位数字 \\
password & string & 是 & 非空字符串 \\
\bottomrule
\end{longtable}

成功响应 \texttt{data}：

\begin{longtable}[]{@{}lll@{}}
\toprule
字段 & 类型 & 约束 \\
\midrule
\endhead
accessToken & string & JWT \\
user.userID & number & 用户 ID \\
user.accountNo & string & 10 位账号 \\
user.displayName & string & 展示昵称 \\
\bottomrule
\end{longtable}

Refresh Token 通过 HttpOnly Cookie 写入。

错误：

\begin{longtable}[]{@{}ll@{}}
\toprule
code & message \\
\midrule
\endhead
VALIDATION\_IDENTIFIER\_INVALID & 账号格式不正确 \\
AUTH\_CREDENTIAL\_INVALID & 账号与密码不匹配 \\
MYSQL\_CONNECTION\_ERROR & 数据库连接失败，登录未完成 \\
\bottomrule
\end{longtable}

\hypertarget{api-auth-003-ux5237ux65b0-token}{%
\subsubsection{API-AUTH-003 刷新
Token}\label{api-auth-003-ux5237ux65b0-token}}

\begin{Shaded}
\begin{Highlighting}[]
\NormalTok{POST /api/v1/auth/refresh}
\end{Highlighting}
\end{Shaded}

请求：

\begin{itemize}
\tightlist
\item
  无请求体。
\item
  必须携带 Refresh Token Cookie。
\end{itemize}

成功响应 \texttt{data}：

\begin{longtable}[]{@{}lll@{}}
\toprule
字段 & 类型 & 约束 \\
\midrule
\endhead
accessToken & string & 新 Access Token \\
\bottomrule
\end{longtable}

错误：

\begin{longtable}[]{@{}ll@{}}
\toprule
code & message \\
\midrule
\endhead
AUTH\_REFRESH\_MISSING & 刷新凭证缺失，请重新登录 \\
AUTH\_REFRESH\_EXPIRED & 登录状态已失效，请重新登录 \\
AUTH\_REFRESH\_REUSE\_DETECTED & 登录状态异常，请重新登录 \\
\bottomrule
\end{longtable}

\hypertarget{api-auth-004-ux9000ux51faux767bux5f55}{%
\subsubsection{API-AUTH-004
退出登录}\label{api-auth-004-ux9000ux51faux767bux5f55}}

\begin{Shaded}
\begin{Highlighting}[]
\NormalTok{POST /api/v1/auth/logout}
\end{Highlighting}
\end{Shaded}

请求：

\begin{itemize}
\tightlist
\item
  必须携带 Refresh Token Cookie。
\end{itemize}

成功响应：

\begin{Shaded}
\begin{Highlighting}[]
\FunctionTok{\{}
  \DataTypeTok{"code"}\FunctionTok{:} \DecValTok{200}\FunctionTok{,}
  \DataTypeTok{"message"}\FunctionTok{:} \StringTok{"退出登录成功"}\FunctionTok{,}
  \DataTypeTok{"data"}\FunctionTok{:} \KeywordTok{null}\FunctionTok{,}
  \DataTypeTok{"traceId"}\FunctionTok{:} \StringTok{"01HZ0000000000000000000000"}
\FunctionTok{\}}
\end{Highlighting}
\end{Shaded}

\hypertarget{api-auth-005-ux5f53ux524dux7528ux6237}{%
\subsubsection{API-AUTH-005
当前用户}\label{api-auth-005-ux5f53ux524dux7528ux6237}}

\begin{Shaded}
\begin{Highlighting}[]
\NormalTok{GET /api/v1/auth/me}
\end{Highlighting}
\end{Shaded}

成功响应 \texttt{data}：

\begin{longtable}[]{@{}lll@{}}
\toprule
字段 & 类型 & 约束 \\
\midrule
\endhead
userID & number & 用户 ID \\
accountNo & string & 10 位账号 \\
displayName & string & 展示昵称 \\
status & string & active \\
\bottomrule
\end{longtable}

\hypertarget{ux5de5ux4f5cux53f0ux4e0eux504fux597dux63a5ux53e3}{%
\subsection{5.
工作台与偏好接口}\label{ux5de5ux4f5cux53f0ux4e0eux504fux597dux63a5ux53e3}}

\hypertarget{api-workbench-001-ux5de5ux4f5cux53f0ux6982ux89c8}{%
\subsubsection{API-WORKBENCH-001
工作台概览}\label{api-workbench-001-ux5de5ux4f5cux53f0ux6982ux89c8}}

\begin{Shaded}
\begin{Highlighting}[]
\NormalTok{GET /api/v1/workbench/overview}
\end{Highlighting}
\end{Shaded}

成功响应 \texttt{data}：

\begin{longtable}[]{@{}lll@{}}
\toprule
字段 & 类型 & 约束 \\
\midrule
\endhead
user.userID & number & 当前用户 ID \\
user.displayName & string & 展示昵称 \\
summary.databaseCount & number & 当前用户数据库数量 \\
summary.tableCount & number & 当前用户表数量 \\
summary.viewCount & number & 当前用户视图数量 \\
summary.indexCount & number & 当前用户索引数量 \\
summary.rowCountEstimated & number & 估算行数 \\
recentOperations & array & 最近操作，按时间倒序 \\
\bottomrule
\end{longtable}

\hypertarget{api-pref-001-ux83b7ux53d6ux504fux597d}{%
\subsubsection{API-PREF-001
获取偏好}\label{api-pref-001-ux83b7ux53d6ux504fux597d}}

\begin{Shaded}
\begin{Highlighting}[]
\NormalTok{GET /api/v1/users/me/preferences}
\end{Highlighting}
\end{Shaded}

成功响应 \texttt{data}：

\begin{longtable}[]{@{}lll@{}}
\toprule
字段 & 类型 & 约束 \\
\midrule
\endhead
themeMode & string & dark、light \\
themeHue & number & 0-359 \\
backgroundPreset & string & particle \\
tablePageSize & number & 1 到服务端最大分页 \\
confirmBatchInsert & boolean & 布尔值 \\
confirmBatchUpdate & boolean & 布尔值 \\
confirmBatchDelete & boolean & 布尔值 \\
confirmCascadeDelete & boolean & 布尔值 \\
\bottomrule
\end{longtable}

\hypertarget{api-pref-002-ux66f4ux65b0ux504fux597d}{%
\subsubsection{API-PREF-002
更新偏好}\label{api-pref-002-ux66f4ux65b0ux504fux597d}}

\begin{Shaded}
\begin{Highlighting}[]
\NormalTok{PATCH /api/v1/users/me/preferences}
\end{Highlighting}
\end{Shaded}

请求体：

\begin{longtable}[]{@{}llll@{}}
\toprule
参数 & 类型 & 必填 & 约束 \\
\midrule
\endhead
themeHue & number & 否 & 0-359 \\
tablePageSize & number & 否 & 1 到服务端最大分页 \\
confirmBatchInsert & boolean & 否 & 布尔值 \\
confirmBatchUpdate & boolean & 否 & 布尔值 \\
confirmBatchDelete & boolean & 否 & 布尔值 \\
confirmCascadeDelete & boolean & 否 & 布尔值 \\
\bottomrule
\end{longtable}

\hypertarget{ux6570ux636eux5e93ux63a5ux53e3}{%
\subsection{6. 数据库接口}\label{ux6570ux636eux5e93ux63a5ux53e3}}

\hypertarget{api-db-001-ux6570ux636eux5e93ux5217ux8868}{%
\subsubsection{API-DB-001
数据库列表}\label{api-db-001-ux6570ux636eux5e93ux5217ux8868}}

\begin{Shaded}
\begin{Highlighting}[]
\NormalTok{GET /api/v1/databases}
\end{Highlighting}
\end{Shaded}

成功响应 \texttt{data.items{[}{]}}：

\begin{longtable}[]{@{}lll@{}}
\toprule
字段 & 类型 & 约束 \\
\midrule
\endhead
id & number & 平台数据库 ID \\
displayName & string & 显示名 \\
charset & string & 字符集 \\
collation & string & 排序规则 \\
tableCount & number & 表数量 \\
viewCount & number & 视图数量 \\
sizeBytes & number & 估算容量，单位 Byte \\
createdAt & string & ISO 时间 \\
updatedAt & string & ISO 时间 \\
\bottomrule
\end{longtable}

\hypertarget{api-db-002-ux521bux5efaux6570ux636eux5e93}{%
\subsubsection{API-DB-002
创建数据库}\label{api-db-002-ux521bux5efaux6570ux636eux5e93}}

\begin{Shaded}
\begin{Highlighting}[]
\NormalTok{POST /api/v1/databases}
\end{Highlighting}
\end{Shaded}

请求体：

\begin{longtable}[]{@{}llll@{}}
\toprule
参数 & 类型 & 必填 & 约束 \\
\midrule
\endhead
displayName & string & 是 & 标识符规则 \\
charset & string & 否 & 默认 utf8mb4 \\
collation & string & 否 & 默认 utf8mb4\_unicode\_ci \\
\bottomrule
\end{longtable}

阶段说明：当前版本仅支持 \texttt{utf8mb4} 与
\texttt{utf8mb4\_unicode\_ci}。

错误：

\begin{longtable}[]{@{}ll@{}}
\toprule
code & message \\
\midrule
\endhead
VALIDATION\_IDENTIFIER\_INVALID & 数据库名称格式不正确 \\
RESOURCE\_CONFLICT & 数据库名称已存在 \\
LIMIT\_EXCEEDED & 数据库数量超过系统限制 \\
MYSQL\_EXECUTION\_ERROR & 数据库创建失败：\{mysql.sqlMessage\} \\
\bottomrule
\end{longtable}

\hypertarget{api-db-003-ux4feeux6539ux6570ux636eux5e93}{%
\subsubsection{API-DB-003
修改数据库}\label{api-db-003-ux4feeux6539ux6570ux636eux5e93}}

\begin{Shaded}
\begin{Highlighting}[]
\NormalTok{PATCH /api/v1/databases/:databaseId}
\end{Highlighting}
\end{Shaded}

路径参数：

\begin{longtable}[]{@{}lll@{}}
\toprule
参数 & 类型 & 约束 \\
\midrule
\endhead
databaseId & number & 正整数 \\
\bottomrule
\end{longtable}

请求体：

\begin{longtable}[]{@{}llll@{}}
\toprule
参数 & 类型 & 必填 & 约束 \\
\midrule
\endhead
displayName & string & 否 & 标识符规则 \\
charset & string & 否 & 当前版本仅支持 utf8mb4 \\
collation & string & 否 & 当前版本仅支持 utf8mb4\_unicode\_ci \\
\bottomrule
\end{longtable}

说明：当前版本只支持修改显示名，字符集和排序规则作为兼容字段校验，不执行物理库变更。

\hypertarget{api-db-004-ux5220ux9664ux6570ux636eux5e93}{%
\subsubsection{API-DB-004
删除数据库}\label{api-db-004-ux5220ux9664ux6570ux636eux5e93}}

\begin{Shaded}
\begin{Highlighting}[]
\NormalTok{DELETE /api/v1/databases/:databaseId}
\end{Highlighting}
\end{Shaded}

请求体：

\begin{longtable}[]{@{}llll@{}}
\toprule
参数 & 类型 & 必填 & 约束 \\
\midrule
\endhead
confirmation.confirmed & boolean & 是 & true \\
confirmation.confirmText & string & 是 & 等于数据库显示名 \\
\bottomrule
\end{longtable}

错误：

\begin{longtable}[]{@{}ll@{}}
\toprule
code & message \\
\midrule
\endhead
CONFIRMATION\_REQUIRED & 删除数据库需要二次确认 \\
CONFIRMATION\_TEXT\_MISMATCH & 数据库名称确认失败 \\
AUTHZ\_DATABASE\_FORBIDDEN & 无权删除该数据库 \\
MYSQL\_EXECUTION\_ERROR & 数据库删除失败：\{mysql.sqlMessage\} \\
\bottomrule
\end{longtable}

\hypertarget{ux8868ux7ed3ux6784ux63a5ux53e3}{%
\subsection{7. 表结构接口}\label{ux8868ux7ed3ux6784ux63a5ux53e3}}

\hypertarget{api-table-001-ux5bf9ux8c61ux5217ux8868}{%
\subsubsection{API-TABLE-001
对象列表}\label{api-table-001-ux5bf9ux8c61ux5217ux8868}}

\begin{Shaded}
\begin{Highlighting}[]
\NormalTok{GET /api/v1/databases/:databaseId/objects}
\end{Highlighting}
\end{Shaded}

查询参数：

\begin{longtable}[]{@{}llll@{}}
\toprule
参数 & 类型 & 必填 & 约束 \\
\midrule
\endhead
objectType & string & 否 & table、view \\
\bottomrule
\end{longtable}

成功响应 \texttt{data.items{[}{]}}：

\begin{longtable}[]{@{}lll@{}}
\toprule
字段 & 类型 & 约束 \\
\midrule
\endhead
objectType & string & table、view \\
name & string & 对象名 \\
rowCountEstimated & number & 估算行数 \\
dataLength & number & 数据大小 \\
indexLength & number & 索引大小 \\
createdAt & string & ISO 时间 \\
\bottomrule
\end{longtable}

\hypertarget{api-table-002-ux521bux5efaux8868}{%
\subsubsection{API-TABLE-002
创建表}\label{api-table-002-ux521bux5efaux8868}}

\begin{Shaded}
\begin{Highlighting}[]
\NormalTok{POST /api/v1/databases/:databaseId/tables}
\end{Highlighting}
\end{Shaded}

请求体：

\begin{longtable}[]{@{}llll@{}}
\toprule
参数 & 类型 & 必填 & 约束 \\
\midrule
\endhead
tableName & string & 是 & 标识符规则 \\
columns & array & 是 & 长度 1 到服务端上限 \\
constraints & array & 否 & 约束定义 \\
indexes & array & 否 & 索引定义 \\
\bottomrule
\end{longtable}

说明：创建表不强制必须包含主键；无主键表允许预览和查询，后续行内更新接口必须按无主键表策略限制。

\texttt{columns{[}{]}}：

\begin{longtable}[]{@{}llll@{}}
\toprule
字段 & 类型 & 必填 & 约束 \\
\midrule
\endhead
name & string & 是 & 标识符规则 \\
type & string & 是 & 支持的数据类型 \\
length & number & 否 & CHAR 使用 1-255，VARCHAR 使用 1-16383 \\
precision & number & 否 & DECIMAL 使用，1-65 \\
scale & number & 否 & DECIMAL 使用，0 到 precision \\
unsigned & boolean & 否 & 数值类型使用 \\
nullable & boolean & 是 & 布尔值 \\
defaultValue & any & 否 & 与字段类型兼容 \\
autoIncrement & boolean & 否 & 整数类型使用 \\
comment & string & 否 & 最大 255 字符 \\
\bottomrule
\end{longtable}

支持的数据类型：

\begin{Shaded}
\begin{Highlighting}[]
\NormalTok{TINYINT, SMALLINT, INT, BIGINT, DECIMAL, FLOAT, DOUBLE,}
\NormalTok{CHAR, VARCHAR, TEXT, MEDIUMTEXT, LONGTEXT,}
\NormalTok{DATE, TIME, DATETIME, TIMESTAMP, BOOLEAN, JSON,}
\NormalTok{BLOB, MEDIUMBLOB, LONGBLOB}
\end{Highlighting}
\end{Shaded}

\texttt{constraints{[}{]}}：

\begin{itemize}
\tightlist
\item
  \texttt{name}：\texttt{string}；必填；标识符规则。
\item
  \texttt{type}：\texttt{string}；必填；取值为
  \texttt{PRIMARY\_KEY}、\texttt{FOREIGN\_KEY}、\texttt{UNIQUE}、\texttt{CHECK}。
\item
  \texttt{columns}：\texttt{array}；必填；字段名数组，字段必须存在。
\item
  \texttt{referencedTable}：\texttt{string}；条件必填；外键使用。
\item
  \texttt{referencedColumns}：\texttt{array}；条件必填；外键使用。
\item
  \texttt{onDelete}：\texttt{string}；条件必填；取值为
  \texttt{RESTRICT}、\texttt{CASCADE}、\texttt{SET\ NULL}、\texttt{NO\ ACTION}。
\item
  \texttt{onUpdate}：\texttt{string}；条件必填；取值为
  \texttt{RESTRICT}、\texttt{CASCADE}、\texttt{SET\ NULL}、\texttt{NO\ ACTION}。
\item
  \texttt{expression}：\texttt{string}；条件必填；\texttt{CHECK}
  使用，后端表达式白名单校验。
\end{itemize}

阶段说明：V4 创建表接口先实现可选 \texttt{PRIMARY\_KEY} 和
\texttt{UNIQUE}；\texttt{FOREIGN\_KEY}、\texttt{CHECK}、字符串有限取值约束、枚举式约束与更完整的约束管理在后续结构管理接口中扩展。

\texttt{indexes{[}{]}}：

\begin{longtable}[]{@{}llll@{}}
\toprule
字段 & 类型 & 必填 & 约束 \\
\midrule
\endhead
name & string & 是 & 标识符规则 \\
unique & boolean & 是 & 布尔值 \\
columns & array & 是 & 索引字段定义 \\
\bottomrule
\end{longtable}

\texttt{indexes{[}{]}.columns{[}{]}}：

\begin{longtable}[]{@{}llll@{}}
\toprule
字段 & 类型 & 必填 & 约束 \\
\midrule
\endhead
name & string & 是 & 字段必须存在 \\
order & string & 是 & ASC、DESC \\
\bottomrule
\end{longtable}

\hypertarget{api-table-003-ux8868ux7ed3ux6784}{%
\subsubsection{API-TABLE-003
表结构}\label{api-table-003-ux8868ux7ed3ux6784}}

\begin{Shaded}
\begin{Highlighting}[]
\NormalTok{GET /api/v1/databases/:databaseId/tables/:tableName/schema}
\end{Highlighting}
\end{Shaded}

成功响应 \texttt{data}：

\begin{longtable}[]{@{}lll@{}}
\toprule
字段 & 类型 & 约束 \\
\midrule
\endhead
tableName & string & 表名 \\
columns & array & 字段列表 \\
constraints & array & 约束列表 \\
indexes & array & 索引列表 \\
\bottomrule
\end{longtable}

\hypertarget{api-table-004-ux4feeux6539ux8868ux7ed3ux6784}{%
\subsubsection{API-TABLE-004
修改表结构}\label{api-table-004-ux4feeux6539ux8868ux7ed3ux6784}}

\begin{Shaded}
\begin{Highlighting}[]
\NormalTok{PATCH /api/v1/databases/:databaseId/tables/:tableName/schema}
\end{Highlighting}
\end{Shaded}

请求体：

\begin{longtable}[]{@{}llll@{}}
\toprule
参数 & 类型 & 必填 & 约束 \\
\midrule
\endhead
operations & array & 是 & 长度 1 到服务端上限 \\
confirmation & object & 条件必填 & 高风险操作必填 \\
\bottomrule
\end{longtable}

\texttt{operations{[}{]}}：

\begin{longtable}[]{@{}ll@{}}
\toprule
action & 必填字段 \\
\midrule
\endhead
ADD\_COLUMN & column \\
MODIFY\_COLUMN & oldName、column \\
RENAME\_COLUMN & oldName、newName \\
DROP\_COLUMN & name \\
ADD\_CONSTRAINT & constraint \\
DROP\_CONSTRAINT & name \\
ADD\_INDEX & index \\
DROP\_INDEX & name \\
\bottomrule
\end{longtable}

错误：

\begin{longtable}[]{@{}ll@{}}
\toprule
code & message \\
\midrule
\endhead
RESOURCE\_NOT\_FOUND & 表不存在 \\
VALIDATION\_IDENTIFIER\_INVALID & 字段名称格式不正确 \\
CONFIRMATION\_REQUIRED & 该结构变更需要二次确认 \\
MYSQL\_EXECUTION\_ERROR & 表结构修改失败：\{mysql.sqlMessage\} \\
\bottomrule
\end{longtable}

\hypertarget{api-table-005-ux5220ux9664ux8868}{%
\subsubsection{API-TABLE-005
删除表}\label{api-table-005-ux5220ux9664ux8868}}

\begin{Shaded}
\begin{Highlighting}[]
\NormalTok{DELETE /api/v1/databases/:databaseId/tables/:tableName}
\end{Highlighting}
\end{Shaded}

请求体：

\begin{longtable}[]{@{}llll@{}}
\toprule
参数 & 类型 & 必填 & 约束 \\
\midrule
\endhead
confirmation.confirmed & boolean & 是 & true \\
confirmation.confirmText & string & 是 & 等于表名 \\
\bottomrule
\end{longtable}

\hypertarget{ux7ea6ux675fux63a5ux53e3}{%
\subsection{8. 约束接口}\label{ux7ea6ux675fux63a5ux53e3}}

\hypertarget{api-constraint-001-ux521bux5efaux7ea6ux675f}{%
\subsubsection{API-CONSTRAINT-001
创建约束}\label{api-constraint-001-ux521bux5efaux7ea6ux675f}}

\begin{Shaded}
\begin{Highlighting}[]
\NormalTok{POST /api/v1/databases/:databaseId/tables/:tableName/constraints}
\end{Highlighting}
\end{Shaded}

请求体同 \texttt{constraints{[}{]}} 单项定义。

\hypertarget{api-constraint-002-ux5220ux9664ux7ea6ux675f}{%
\subsubsection{API-CONSTRAINT-002
删除约束}\label{api-constraint-002-ux5220ux9664ux7ea6ux675f}}

\begin{Shaded}
\begin{Highlighting}[]
\NormalTok{DELETE /api/v1/databases/:databaseId/tables/:tableName/constraints/:constraintName}
\end{Highlighting}
\end{Shaded}

请求体：

\begin{longtable}[]{@{}llll@{}}
\toprule
参数 & 类型 & 必填 & 约束 \\
\midrule
\endhead
confirmation.confirmed & boolean & 条件必填 & 高风险约束必填 true \\
confirmation.confirmText & string & 条件必填 & 等于约束名 \\
\bottomrule
\end{longtable}

\hypertarget{ux7d22ux5f15ux63a5ux53e3}{%
\subsection{9. 索引接口}\label{ux7d22ux5f15ux63a5ux53e3}}

\hypertarget{api-index-001-ux521bux5efaux7d22ux5f15}{%
\subsubsection{API-INDEX-001
创建索引}\label{api-index-001-ux521bux5efaux7d22ux5f15}}

\begin{Shaded}
\begin{Highlighting}[]
\NormalTok{POST /api/v1/databases/:databaseId/tables/:tableName/indexes}
\end{Highlighting}
\end{Shaded}

请求体同 \texttt{indexes{[}{]}} 单项定义。

\hypertarget{api-index-002-ux5220ux9664ux7d22ux5f15}{%
\subsubsection{API-INDEX-002
删除索引}\label{api-index-002-ux5220ux9664ux7d22ux5f15}}

\begin{Shaded}
\begin{Highlighting}[]
\NormalTok{DELETE /api/v1/databases/:databaseId/tables/:tableName/indexes/:indexName}
\end{Highlighting}
\end{Shaded}

成功响应 \texttt{data} 固定为 null。

\hypertarget{ux89c6ux56feux63a5ux53e3}{%
\subsection{10. 视图接口}\label{ux89c6ux56feux63a5ux53e3}}

\hypertarget{api-view-001-ux89c6ux56feux5217ux8868}{%
\subsubsection{API-VIEW-001
视图列表}\label{api-view-001-ux89c6ux56feux5217ux8868}}

\begin{Shaded}
\begin{Highlighting}[]
\NormalTok{GET /api/v1/databases/:databaseId/views}
\end{Highlighting}
\end{Shaded}

成功响应 \texttt{data.items{[}{]}}：

\begin{longtable}[]{@{}lll@{}}
\toprule
字段 & 类型 & 约束 \\
\midrule
\endhead
viewName & string & 视图名 \\
updatable & boolean & 是否可更新 \\
createdAt & string & ISO 时间 \\
\bottomrule
\end{longtable}

\hypertarget{api-view-002-ux521bux5efaux89c6ux56fe}{%
\subsubsection{API-VIEW-002
创建视图}\label{api-view-002-ux521bux5efaux89c6ux56fe}}

\begin{Shaded}
\begin{Highlighting}[]
\NormalTok{POST /api/v1/databases/:databaseId/views}
\end{Highlighting}
\end{Shaded}

请求体：

\begin{longtable}[]{@{}llll@{}}
\toprule
参数 & 类型 & 必填 & 约束 \\
\midrule
\endhead
viewName & string & 是 & 标识符规则 \\
selectDefinition & object & 是 & 同 API-QUERY-001 \\
\bottomrule
\end{longtable}

\hypertarget{api-view-003-ux4feeux6539ux89c6ux56fe}{%
\subsubsection{API-VIEW-003
修改视图}\label{api-view-003-ux4feeux6539ux89c6ux56fe}}

\begin{Shaded}
\begin{Highlighting}[]
\NormalTok{PATCH /api/v1/databases/:databaseId/views/:viewName}
\end{Highlighting}
\end{Shaded}

请求体：

\begin{longtable}[]{@{}llll@{}}
\toprule
参数 & 类型 & 必填 & 约束 \\
\midrule
\endhead
selectDefinition & object & 是 & 同 API-QUERY-001 \\
\bottomrule
\end{longtable}

\hypertarget{api-view-004-ux5220ux9664ux89c6ux56fe}{%
\subsubsection{API-VIEW-004
删除视图}\label{api-view-004-ux5220ux9664ux89c6ux56fe}}

\begin{Shaded}
\begin{Highlighting}[]
\NormalTok{DELETE /api/v1/databases/:databaseId/views/:viewName}
\end{Highlighting}
\end{Shaded}

请求体：

\begin{longtable}[]{@{}llll@{}}
\toprule
参数 & 类型 & 必填 & 约束 \\
\midrule
\endhead
confirmation.confirmed & boolean & 是 & true \\
confirmation.confirmText & string & 是 & 等于视图名 \\
\bottomrule
\end{longtable}

\hypertarget{ux8868ux6570ux636eux63a5ux53e3}{%
\subsection{11. 表数据接口}\label{ux8868ux6570ux636eux63a5ux53e3}}

\hypertarget{api-row-001-ux9884ux89c8ux8868ux6570ux636e}{%
\subsubsection{API-ROW-001
预览表数据}\label{api-row-001-ux9884ux89c8ux8868ux6570ux636e}}

\begin{Shaded}
\begin{Highlighting}[]
\NormalTok{GET /api/v1/databases/:databaseId/tables/:tableName/preview}
\end{Highlighting}
\end{Shaded}

查询参数：

\begin{longtable}[]{@{}llll@{}}
\toprule
参数 & 类型 & 必填 & 约束 \\
\midrule
\endhead
limit & number & 否 & 1 到服务端分页上限，默认 20 \\
offset & number & 否 & 0 到查询行数上限，默认 0 \\
filters & string & 否 & JSON 字符串，字段筛选数组 \\
\bottomrule
\end{longtable}

\texttt{filters{[}{]}}：

\begin{longtable}[]{@{}llll@{}}
\toprule
字段 & 类型 & 必填 & 约束 \\
\midrule
\endhead
column & string & 是 & 必须是目标表真实字段 \\
value & string & 是 & 非空 \\
mode & string & 否 & \texttt{equals} 或 \texttt{contains}，默认
\texttt{contains} \\
\bottomrule
\end{longtable}

成功响应 \texttt{data}：

\begin{longtable}[]{@{}lll@{}}
\toprule
字段 & 类型 & 约束 \\
\midrule
\endhead
tableName & string & 表名 \\
columns & array & 字段结构，格式同 \texttt{API-TABLE-003} \\
rows & array & 当前页行数据 \\
total & number & 当前筛选条件下总行数 \\
limit & number & 本次请求行数上限 \\
offset & number & 本次请求偏移 \\
hasMore & boolean & 是否还有更多行 \\
facets & object & 有限离散字段候选值，用于下拉筛选 \\
\bottomrule
\end{longtable}

说明：

\begin{itemize}
\tightlist
\item
  V5 阶段使用预览接口承载分页浏览、有限值筛选和包含搜索。
\item
  留空筛选条件表示全选；前端按 Enter 或点击刷新后才重新请求。
\item
  后续正式查询构造器会在 \texttt{API-QUERY}
  中支持排序、JOIN、聚合、子查询和更复杂条件。
\end{itemize}

\hypertarget{api-row-002-ux65b0ux589eux5355ux884cux6570ux636e}{%
\subsubsection{API-ROW-002
新增单行数据}\label{api-row-002-ux65b0ux589eux5355ux884cux6570ux636e}}

\begin{Shaded}
\begin{Highlighting}[]
\NormalTok{POST /api/v1/databases/:databaseId/tables/:tableName/rows}
\end{Highlighting}
\end{Shaded}

请求体：

\begin{longtable}[]{@{}llll@{}}
\toprule
参数 & 类型 & 必填 & 约束 \\
\midrule
\endhead
row & object & 是 & 字段名到字段值的对象，至少包含一个字段 \\
\bottomrule
\end{longtable}

规则：

\begin{itemize}
\tightlist
\item
  字段名必须存在于目标表。
\item
  \texttt{auto\_increment} 字段可以不提交，由 MySQL 自动生成。
\item
  前端新增行时空白字段默认不提交，避免误写入 \texttt{NULL}。
\item
  后端按表结构校验数字、布尔、JSON 和非空字段。
\item
  单表容量达到限制时拒绝新增。
\end{itemize}

成功响应 \texttt{data}：

\begin{longtable}[]{@{}lll@{}}
\toprule
字段 & 类型 & 约束 \\
\midrule
\endhead
tableName & string & 表名 \\
row & object 或 null & 可通过主键回读时返回新增后的行 \\
primaryKey & object 或 null & 有主键时返回主键值，无主键时为 null \\
\bottomrule
\end{longtable}

\hypertarget{api-row-003-ux4feeux6539ux5355ux884cux6570ux636e}{%
\subsubsection{API-ROW-003
修改单行数据}\label{api-row-003-ux4feeux6539ux5355ux884cux6570ux636e}}

\begin{Shaded}
\begin{Highlighting}[]
\NormalTok{PATCH /api/v1/databases/:databaseId/tables/:tableName/rows}
\end{Highlighting}
\end{Shaded}

请求体：

\begin{longtable}[]{@{}llll@{}}
\toprule
参数 & 类型 & 必填 & 约束 \\
\midrule
\endhead
primaryKey & object & 是 & 完整主键字段和值 \\
set & object & 是 & 字段名到新值的对象 \\
\bottomrule
\end{longtable}

行内编辑提交规则：

\begin{itemize}
\tightlist
\item
  用户在前端表格中直接修改单元格时，前端只保存本地草稿，不立即调用该接口。
\item
  用户点击该行旁边``更新/提交''按钮后，前端调用该接口。
\item
  后端使用 \texttt{primaryKey} 定位唯一行。
\item
  \texttt{set} 只提交被修改字段，不提交整行数据。
\item
  \texttt{primaryKey} 字段和 \texttt{auto\_increment}
  字段不允许通过行内编辑修改。
\item
  用户点击``取消''时前端丢弃草稿，不调用后端。
\item
  无主键表在 V5 阶段禁止行内编辑。
\item
  单表容量达到限制时拒绝更新。
\end{itemize}

行内编辑请求示例：

\begin{Shaded}
\begin{Highlighting}[]
\FunctionTok{\{}
  \DataTypeTok{"primaryKey"}\FunctionTok{:} \FunctionTok{\{}
    \DataTypeTok{"id"}\FunctionTok{:} \DecValTok{1}
  \FunctionTok{\},}
  \DataTypeTok{"set"}\FunctionTok{:} \FunctionTok{\{}
    \DataTypeTok{"name"}\FunctionTok{:} \StringTok{"Alice"}\FunctionTok{,}
    \DataTypeTok{"age"}\FunctionTok{:} \DecValTok{18}
  \FunctionTok{\}}
\FunctionTok{\}}
\end{Highlighting}
\end{Shaded}

错误：

\begin{longtable}[]{@{}ll@{}}
\toprule
code & message \\
\midrule
\endhead
PRIMARY\_KEY\_REQUIRED & 当前表没有主键，暂不支持行级修改或删除 \\
ROW\_UPDATE\_NOT\_UNIQUE & 更新目标不存在或不唯一，请刷新后重试 \\
VALIDATION\_FIELD\_INVALID & 字段不存在、类型不合法或字段不允许修改 \\
MYSQL\_EXECUTION\_ERROR & 更新行失败：\{mysql.sqlMessage\} \\
\bottomrule
\end{longtable}

\hypertarget{api-row-004-ux5220ux9664ux5355ux884cux6570ux636e}{%
\subsubsection{API-ROW-004
删除单行数据}\label{api-row-004-ux5220ux9664ux5355ux884cux6570ux636e}}

\begin{Shaded}
\begin{Highlighting}[]
\NormalTok{DELETE /api/v1/databases/:databaseId/tables/:tableName/rows}
\end{Highlighting}
\end{Shaded}

请求体：

\begin{longtable}[]{@{}llll@{}}
\toprule
参数 & 类型 & 必填 & 约束 \\
\midrule
\endhead
primaryKey & object & 是 & 完整主键字段和值 \\
\bottomrule
\end{longtable}

规则：

\begin{itemize}
\tightlist
\item
  后端使用 \texttt{primaryKey} 定位唯一行。
\item
  无主键表在 V5 阶段禁止单行删除。
\item
  前端删除前使用玻璃风格确认窗口展示主键信息。
\end{itemize}

错误：

\begin{longtable}[]{@{}ll@{}}
\toprule
code & message \\
\midrule
\endhead
PRIMARY\_KEY\_REQUIRED & 当前表没有主键，暂不支持行级修改或删除 \\
ROW\_DELETE\_NOT\_UNIQUE & 删除目标不存在或不唯一，请刷新后重试 \\
MYSQL\_EXECUTION\_ERROR & 删除行失败：\{mysql.sqlMessage\} \\
\bottomrule
\end{longtable}

\hypertarget{api-row-005-ux6279ux91cfux6570ux636eux64cdux4f5cux89c4ux5212}{%
\subsubsection{API-ROW-005
批量数据操作（规划）}\label{api-row-005-ux6279ux91cfux6570ux636eux64cdux4f5cux89c4ux5212}}

批量新增、批量条件修改和批量条件删除作为后续增强接口落地。规划要求如下：

\begin{itemize}
\tightlist
\item
  批量新增必须限制请求行数，并在失败时返回失败行号、失败字段、MySQL
  错误码和数据库错误消息。
\item
  批量修改和批量删除必须使用 GUI 条件对象，不允许无条件操作。
\item
  影响行数超过阈值时必须返回
  \texttt{CONFIRMATION\_REQUIRED}，由前端二次确认后继续。
\item
  涉及级联删除时必须展示级联影响范围。
\end{itemize}

\hypertarget{gui-ux67e5ux8be2ux63a5ux53e3}{%
\subsection{12. GUI 查询接口}\label{gui-ux67e5ux8be2ux63a5ux53e3}}

\hypertarget{api-query-001-ux7ed3ux6784ux5316-select}{%
\subsubsection{API-QUERY-001 结构化
SELECT}\label{api-query-001-ux7ed3ux6784ux5316-select}}

\begin{Shaded}
\begin{Highlighting}[]
\NormalTok{POST /api/v1/databases/:databaseId/query/select}
\end{Highlighting}
\end{Shaded}

请求体：

\begin{longtable}[]{@{}llll@{}}
\toprule
参数 & 类型 & 必填 & 约束 \\
\midrule
\endhead
from & object & 是 & 主查询来源，支持表、视图、派生表子查询 \\
joins & array & 否 & 连接定义 \\
fields & array & 是 & 查询字段，长度 1 到服务端上限 \\
filters & array & 否 & 筛选条件 \\
groups & array & 否 & 分组字段 \\
having & array & 否 & 聚合筛选 \\
sorts & array & 否 & 排序 \\
page & number & 否 & 最小 1 \\
pageSize & number & 否 & 服务端最大值限制 \\
\bottomrule
\end{longtable}

\texttt{from} 和 \texttt{joins{[}{]}.source} 使用统一来源对象：

\begin{longtable}[]{@{}llll@{}}
\toprule
字段 & 类型 & 必填 & 约束 \\
\midrule
\endhead
type & string & 是 & table、view、subquery \\
table & string & 条件必填 & type=table 时必填，表名存在 \\
view & string & 条件必填 & type=view 时必填，视图名存在 \\
query & object & 条件必填 & type=subquery 时必填，结构同 API-QUERY-001
请求体 \\
alias & string & 否 & type=subquery 时必填，标识符规则 \\
\bottomrule
\end{longtable}

\texttt{joins{[}{]}}：

\begin{longtable}[]{@{}llll@{}}
\toprule
字段 & 类型 & 必填 & 约束 \\
\midrule
\endhead
type & string & 是 & INNER、LEFT、RIGHT \\
source & object & 是 & 来源对象，支持表、视图、子查询 \\
alias & string & 是 & 标识符规则 \\
on & array & 是 & 连接条件 \\
\bottomrule
\end{longtable}

\texttt{fields{[}{]}}：

\begin{longtable}[]{@{}llll@{}}
\toprule
字段 & 类型 & 必填 & 约束 \\
\midrule
\endhead
tableAlias & string & 否 & 已声明别名 \\
name & string & 是 & 字段存在 \\
alias & string & 否 & 标识符规则 \\
aggregate & string & 否 & COUNT、SUM、AVG、MIN、MAX \\
\bottomrule
\end{longtable}

\texttt{filters{[}{]}} 和 \texttt{having{[}{]}}：

\begin{longtable}[]{@{}
  >{\raggedright\arraybackslash}p{(\columnwidth - 6\tabcolsep) * \real{0.25}}
  >{\raggedright\arraybackslash}p{(\columnwidth - 6\tabcolsep) * \real{0.25}}
  >{\raggedright\arraybackslash}p{(\columnwidth - 6\tabcolsep) * \real{0.25}}
  >{\raggedright\arraybackslash}p{(\columnwidth - 6\tabcolsep) * \real{0.25}}@{}}
\toprule
字段 & 类型 & 必填 & 约束 \\
\midrule
\endhead
field & string & 是 & 字段存在 \\
operator & string & 是 &
eq、ne、gt、gte、lt、lte、like、in、notIn、between、isNull、isNotNull、exists、notExists \\
value & any & 条件必填 & 与字段类型兼容 \\
subquery & object & 条件必填 &
in、notIn、exists、notExists、标量比较时使用 \\
logic & string & 否 & AND、OR \\
\bottomrule
\end{longtable}

子查询规则：

\begin{itemize}
\tightlist
\item
  子查询本身也是结构化查询 AST。
\item
  支持派生表子查询、\texttt{IN} / \texttt{NOT\ IN}
  子查询、\texttt{EXISTS} / \texttt{NOT\ EXISTS} 子查询和标量子查询。
\item
  最大子查询嵌套深度为 3。
\item
  子查询和主查询合计 JOIN 数不得超过 8。
\item
  所有子查询仍然只能生成 \texttt{SELECT} SQL。
\item
  子查询引用字段必须能在当前查询作用域中解析。
\end{itemize}

结构化子查询示例：

\begin{Shaded}
\begin{Highlighting}[]
\FunctionTok{\{}
  \DataTypeTok{"from"}\FunctionTok{:} \FunctionTok{\{}
    \DataTypeTok{"type"}\FunctionTok{:} \StringTok{"table"}\FunctionTok{,}
    \DataTypeTok{"table"}\FunctionTok{:} \StringTok{"students"}\FunctionTok{,}
    \DataTypeTok{"alias"}\FunctionTok{:} \StringTok{"s"}
  \FunctionTok{\},}
  \DataTypeTok{"fields"}\FunctionTok{:} \OtherTok{[}
    \FunctionTok{\{}
      \DataTypeTok{"tableAlias"}\FunctionTok{:} \StringTok{"s"}\FunctionTok{,}
      \DataTypeTok{"name"}\FunctionTok{:} \StringTok{"name"}
    \FunctionTok{\}}
  \OtherTok{]}\FunctionTok{,}
  \DataTypeTok{"filters"}\FunctionTok{:} \OtherTok{[}
    \FunctionTok{\{}
      \DataTypeTok{"field"}\FunctionTok{:} \StringTok{"s.id"}\FunctionTok{,}
      \DataTypeTok{"operator"}\FunctionTok{:} \StringTok{"in"}\FunctionTok{,}
      \DataTypeTok{"subquery"}\FunctionTok{:} \FunctionTok{\{}
        \DataTypeTok{"from"}\FunctionTok{:} \FunctionTok{\{}
          \DataTypeTok{"type"}\FunctionTok{:} \StringTok{"table"}\FunctionTok{,}
          \DataTypeTok{"table"}\FunctionTok{:} \StringTok{"scores"}\FunctionTok{,}
          \DataTypeTok{"alias"}\FunctionTok{:} \StringTok{"sc"}
        \FunctionTok{\},}
        \DataTypeTok{"fields"}\FunctionTok{:} \OtherTok{[}
          \FunctionTok{\{}
            \DataTypeTok{"tableAlias"}\FunctionTok{:} \StringTok{"sc"}\FunctionTok{,}
            \DataTypeTok{"name"}\FunctionTok{:} \StringTok{"student\_id"}
          \FunctionTok{\}}
        \OtherTok{]}\FunctionTok{,}
        \DataTypeTok{"filters"}\FunctionTok{:} \OtherTok{[}
          \FunctionTok{\{}
            \DataTypeTok{"field"}\FunctionTok{:} \StringTok{"sc.score"}\FunctionTok{,}
            \DataTypeTok{"operator"}\FunctionTok{:} \StringTok{"gte"}\FunctionTok{,}
            \DataTypeTok{"value"}\FunctionTok{:} \DecValTok{60}
          \FunctionTok{\}}
        \OtherTok{]}
      \FunctionTok{\}}
    \FunctionTok{\}}
  \OtherTok{]}\FunctionTok{,}
  \DataTypeTok{"page"}\FunctionTok{:} \DecValTok{1}\FunctionTok{,}
  \DataTypeTok{"pageSize"}\FunctionTok{:} \DecValTok{20}
\FunctionTok{\}}
\end{Highlighting}
\end{Shaded}

\hypertarget{ux4e8bux52a1ux63a5ux53e3}{%
\subsection{13. 事务接口}\label{ux4e8bux52a1ux63a5ux53e3}}

\hypertarget{api-tx-001-ux5f00ux542fux4e8bux52a1}{%
\subsubsection{API-TX-001
开启事务}\label{api-tx-001-ux5f00ux542fux4e8bux52a1}}

\begin{Shaded}
\begin{Highlighting}[]
\NormalTok{POST /api/v1/databases/:databaseId/transactions}
\end{Highlighting}
\end{Shaded}

请求体：

\begin{itemize}
\tightlist
\item
  \texttt{isolationLevel}：\texttt{string}；必填；取值为
  \texttt{READ\ COMMITTED}、\texttt{REPEATABLE\ READ}、\texttt{SERIALIZABLE}。
\end{itemize}

成功响应 \texttt{data}：

\begin{longtable}[]{@{}lll@{}}
\toprule
字段 & 类型 & 约束 \\
\midrule
\endhead
transactionId & string & 事务 ID \\
status & string & active \\
expiresAt & string & ISO 时间 \\
\bottomrule
\end{longtable}

\hypertarget{api-tx-002-ux63d0ux4ea4ux4e8bux52a1}{%
\subsubsection{API-TX-002
提交事务}\label{api-tx-002-ux63d0ux4ea4ux4e8bux52a1}}

\begin{Shaded}
\begin{Highlighting}[]
\NormalTok{POST /api/v1/transactions/:transactionId/commit}
\end{Highlighting}
\end{Shaded}

\hypertarget{api-tx-003-ux56deux6edaux4e8bux52a1}{%
\subsubsection{API-TX-003
回滚事务}\label{api-tx-003-ux56deux6edaux4e8bux52a1}}

\begin{Shaded}
\begin{Highlighting}[]
\NormalTok{POST /api/v1/transactions/:transactionId/rollback}
\end{Highlighting}
\end{Shaded}

错误：

\begin{longtable}[]{@{}ll@{}}
\toprule
code & message \\
\midrule
\endhead
RESOURCE\_NOT\_FOUND & 事务不存在 \\
TRANSACTION\_EXPIRED & 事务已超时并回滚 \\
AUTHZ\_TRANSACTION\_FORBIDDEN & 无权操作该事务 \\
\bottomrule
\end{longtable}

\hypertarget{ux7edfux8ba1ux63a5ux53e3}{%
\subsection{14. 统计接口}\label{ux7edfux8ba1ux63a5ux53e3}}

\hypertarget{api-stats-001-ux7528ux6237ux7edfux8ba1}{%
\subsubsection{API-STATS-001
用户统计}\label{api-stats-001-ux7528ux6237ux7edfux8ba1}}

\begin{Shaded}
\begin{Highlighting}[]
\NormalTok{GET /api/v1/stats/user}
\end{Highlighting}
\end{Shaded}

返回数据库数量、表数量、视图数量、索引数量、估算行数、操作趋势。

\hypertarget{api-stats-002-ux6570ux636eux5e93ux7edfux8ba1}{%
\subsubsection{API-STATS-002
数据库统计}\label{api-stats-002-ux6570ux636eux5e93ux7edfux8ba1}}

\begin{Shaded}
\begin{Highlighting}[]
\NormalTok{GET /api/v1/stats/databases/:databaseId}
\end{Highlighting}
\end{Shaded}

返回表数量、视图数量、索引数量、数据大小、索引大小、字段类型分布。

\hypertarget{api-stats-003-ux8868ux7edfux8ba1}{%
\subsubsection{API-STATS-003
表统计}\label{api-stats-003-ux8868ux7edfux8ba1}}

\begin{Shaded}
\begin{Highlighting}[]
\NormalTok{GET /api/v1/stats/databases/:databaseId/tables/:tableName}
\end{Highlighting}
\end{Shaded}

返回字段数量、主键信息、外键信息、索引信息、空值统计、数值聚合统计。

\hypertarget{ux5ba1ux8ba1ux63a5ux53e3}{%
\subsection{15. 审计接口}\label{ux5ba1ux8ba1ux63a5ux53e3}}

\hypertarget{api-audit-001-ux64cdux4f5cux65e5ux5fd7ux5217ux8868}{%
\subsubsection{API-AUDIT-001
操作日志列表}\label{api-audit-001-ux64cdux4f5cux65e5ux5fd7ux5217ux8868}}

\begin{Shaded}
\begin{Highlighting}[]
\NormalTok{GET /api/v1/audit{-}logs}
\end{Highlighting}
\end{Shaded}

查询参数：

\begin{longtable}[]{@{}llll@{}}
\toprule
参数 & 类型 & 必填 & 约束 \\
\midrule
\endhead
page & number & 否 & 最小 1 \\
pageSize & number & 否 & 服务端最大值限制 \\
actionType & string & 否 & 审计动作枚举 \\
objectType & string & 否 & 对象类型枚举 \\
databaseId & number & 否 & 正整数 \\
traceId & string & 否 & 追踪 ID \\
startDate & string & 否 & ISO 时间 \\
endDate & string & 否 & ISO 时间 \\
\bottomrule
\end{longtable}

成功响应 \texttt{data.items{[}{]}}：

\begin{longtable}[]{@{}lll@{}}
\toprule
字段 & 类型 & 约束 \\
\midrule
\endhead
id & number & 日志 ID \\
traceId & string & 追踪 ID \\
actionType & string & 操作类型 \\
objectType & string & 对象类型 \\
objectName & string & 对象名 \\
summary & string & 摘要 \\
success & boolean & 成功状态 \\
errorCode & string & 错误码 \\
createdAt & string & ISO 时间 \\
\bottomrule
\end{longtable}

\end{document}
